% Gerado por scripts/migrate_markdown_to_latex.py.
% Fonte migrada: docs/checklist_pre_benchmark.md


\chapter{Checklist pré-benchmark}





Atualizado em 14/09/2026. 'Benchmark principal' significa a coleta destinada a sustentar conclusões de desempenho no hardware-alvo. Smoke tests em CPU são permitidos para validar a infraestrutura, desde que sejam rotulados como tal. A validação conceitual, o benchmark GPU v1, a robustez sintética v3 e a bateria física v3 já foram concluídos; este checklist agora controla a coleta final dos modelos OPT.



\section{Obrigatório antes do benchmark principal}


\begin{itemize}
\item \textbf{[OBRIGATÓRIO ANTES DO BENCHMARK]} Arquivar a reunião e separar pedidos
explícitos, decisões, dúvidas e interpretações.
\item \textbf{[OBRIGATÓRIO ANTES DO BENCHMARK]} Validar a atenção manual contra
referência matemática, PyTorch e TensorFlow/Keras com as mesmas entradas.
\item \textbf{[OBRIGATÓRIO ANTES DO BENCHMARK]} Registrar tabela, tolerâncias,
versões e limites da validação cruzada.
\item \textbf{[OBRIGATÓRIO ANTES DO BENCHMARK]} Fixar o recorte: projeção densa como
operação prevista no plano e self-attention como foco principal; bloco
Transformer completo fora da coleta principal.
\item \textbf{[OBRIGATÓRIO ANTES DO BENCHMARK]} Fixar cenários principais: baseline,
pruning por magnitude de 50\% e quantização INT8 simulada.
\item \textbf{[OBRIGATÓRIO ANTES DO BENCHMARK]} Implementar runner das operações
isoladas sem alterar o schema CSV existente.
\item \textbf{[OBRIGATÓRIO ANTES DO BENCHMARK]} Versionar configurações distintas de
smoke e coleta principal.
\item \textbf{[OBRIGATÓRIO ANTES DO BENCHMARK]} Codificar e testar o contrato formal:
seed 42,
lote 1, grade \texttt{L=\{64,128,256\}}, \texttt{D=\{128,256,512\}}, 20 warm-ups, 50 medições e
duas execuções independentes. A execução da grade continua condicionada à GPU.
\item \textbf{[OBRIGATÓRIO ANTES DO BENCHMARK]} Registrar Python, bibliotecas,
sistema, dispositivo, parâmetros, ordem dos cenários e método de percentil.
\item \textbf{[OBRIGATÓRIO ANTES DO BENCHMARK]} Persistir CSV, metadados, manifesto,
log e hashes de entradas/saídas em pasta versionável.
\item \textbf{[OBRIGATÓRIO ANTES DO BENCHMARK]} Rejeitar shapes, dtypes, métricas ou
resultados inválidos, \texttt{NaN} e infinitos.
\item \textbf{[OBRIGATÓRIO ANTES DO BENCHMARK]} Corrigir e documentar a inicialização
dos pesos-base: o diagnóstico v1 mostrou que \texttt{N(0,1)} faz a escala da saída
crescer com a dimensão e prejudica a comparação de MSE/MAE.
\item \textbf{[OBRIGATÓRIO ANTES DO BENCHMARK]} Fazer smoke test ponta a ponta,
incluindo análise, gráficos e reexecução determinística.
\item \textbf{[OBRIGATÓRIO ANTES DO BENCHMARK]} Confirmar que o dispositivo da coleta
principal é CUDA e registrar a RTX 4070 Ti Super. A GPU e o runtime NVIDIA no
Docker WSL2 foram confirmados pelo probe; smoke e bateria física foram executados
com a GPU ociosa.
\end{itemize}


\section{Recomendado antes do benchmark principal}


\begin{itemize}
\item \textbf{[RECOMENDADO ANTES DO BENCHMARK]} Variar deterministicamente a ordem
dos cenários entre execuções para reduzir viés de ordem.
\item \textbf{[RECOMENDADO ANTES DO BENCHMARK]} Registrar snapshot exato do ambiente
instalado, sem transformar o arquivo de requisitos mínimos em um lock
específico desta máquina.
\item \textbf{[RECOMENDADO ANTES DO BENCHMARK]} Fazer sanity checks de crescimento de
FLOPs e shapes antes da grade completa.
\item \textbf{[RECOMENDADO ANTES DO BENCHMARK]} Consolidar as duas execuções com
média, desvio-padrão e resultados individuais preservados.
\item \textbf{[RECOMENDADO ANTES DO BENCHMARK]} Confirmar datas institucionais e
corrigir, no canal apropriado, a inconsistência 2026/2027 do cronograma.
\end{itemize}


\section{Pode ser feito depois dos primeiros benchmarks}


\begin{itemize}
\item \textbf{[AMPLIAÇÃO V3 EM COLETA]} Inferência com OPT-125M, OPT-350M e
OPT-1.3B pré-treinados, com pesos e dataset armazenados e smoke funcional aprovado.
\item \textbf{[PODE SER FEITO DEPOIS]} Benchmark de bloco Transformer completo.
\item \textbf{[AMPLIAÇÃO V3 CONCLUÍDA NAS OPERAÇÕES]} Cenários TorchAO INT8 e
pruning 2:4 confirmados pelo profiler na bateria física de 1.650 registros.
\item \textbf{[AMPLIAÇÃO V3 CONCLUÍDA]} Cinco seeds, três repetições, múltiplos
lotes, cabeças, larguras, sequências e distribuições foram avaliados na bateria
sintética de 2.340 registros.
\item \textbf{[PODE SER FEITO DEPOIS]} Comparação direta de desempenho entre
TensorFlow e PyTorch; TensorFlow é apenas referência de correção nesta etapa.
\end{itemize}


\section{Regra de liberação}







O benchmark físico de operações foi liberado e concluído após a aprovação da
imagem fixada, do probe e do smoke. A coleta OPT final mantém o mesmo portão: jogo
fechado, uso médio de cinco amostras abaixo de 10\%, até 2.048~MiB de VRAM ocupada
e temperatura abaixo de 65~graus Celsius. Em WDDM, leitura residual só é aceita
com processos abaixo de 10\%, potência até 35~W e clock até 300~MHz, mantendo os
limites de memória e temperatura. O pico também é registrado no diagnóstico.
Falhas e técnicas incompatíveis devem ser preservadas, sem alegação de desempenho
de hardware quando o kernel correspondente não for confirmado.
